
\documentclass{article}
\usepackage{colortbl}
\usepackage{makecell}
\usepackage{multirow}
\usepackage{supertabular}

\begin{document}

\newcounter{utterance}

\centering \large Interaction Transcript for game `eqbench', experiment `eq\_bench', episode 158 with reflexion{-}llama{-}3.1{-}8b{-}instruct.
\vspace{24pt}

{ \footnotesize  \setcounter{utterance}{1}
\setlength{\tabcolsep}{0pt}
\begin{supertabular}{c@{$\;$}|p{.15\linewidth}@{}p{.15\linewidth}p{.15\linewidth}p{.15\linewidth}p{.15\linewidth}p{.15\linewidth}}
    \# & \multicolumn{2}{c}{Player} && \multicolumn{2}{c}{Game Master} \\
    \hline

    \theutterance \stepcounter{utterance}  
    & & & \multicolumn{4}{p{0.6\linewidth}}{
        \cellcolor[rgb]{0.9,0.9,0.9}{
            \makecell[{{p{\linewidth}}}]{
                \texttt{\tiny{[P1$\langle$GM]}}
                \texttt{You have attempted this same task before and did not succeed. These are the plans you wrote after those attempts.} \\
\texttt{Trial \#1: New plan: } \\
\\ 
\texttt{Upon re{-}reading the dialogue, I realized that I focused too much on the surface{-}level banter between Remy and Eilis, without delving deeper into their underlying emotions and motivations. I should have paid more attention to the subtle cues and power dynamics at play in their conversation.} \\
\\ 
\texttt{To improve my analysis, I will focus on identifying the emotional triggers and vulnerabilities that Remy and Eilis expose throughout the dialogue. I will also consider the context of their relationship and the stakes of their wager, which may have influenced their emotional responses.} \\
\\ 
\texttt{Specifically, I will:} \\
\\ 
\texttt{1. Identify the moments where Remy and Eilis reveal their emotional vulnerabilities, such as Eilis's fear of failure and Remy's romanticism.} \\
\texttt{2. Analyze the power dynamics at play in their conversation, including the ways in which they use language to assert dominance or challenge each other.} \\
\texttt{3. Consider the context of their relationship and the implications of their wager, which may have influenced their emotional responses.} \\
\\ 
\texttt{By taking a more nuanced and contextual approach to analyzing the dialogue, I hope to improve my understanding of Remy's emotional state at the end of the conversation.} \\
\\ 
\texttt{Trial \#2: New plan: } \\
\\ 
\texttt{Upon re{-}reading the dialogue, I realized that I focused too much on the surface{-}level banter between Remy and Eilis, without delving deeper into their underlying emotions and motivations. I should have paid more attention to the subtle cues and power dynamics at play in their conversation.} \\
\\ 
\texttt{To improve my analysis, I will focus on identifying the emotional triggers and vulnerabilities that Remy and Eilis expose throughout the dialogue. I will also consider the context of their relationship and the stakes of their wager, which may have influenced their emotional responses.} \\
\\ 
\texttt{Specifically, I will:} \\
\\ 
\texttt{1. Identify the moments where Remy and Eilis reveal their emotional vulnerabilities, such as Eilis's fear of failure and Remy's romanticism.} \\
\texttt{2. Analyze the power dynamics at play in their conversation, including the ways in which they use language to assert dominance or challenge each other.} \\
\texttt{3. Consider the context of their relationship and the implications of their wager, which may have influenced their emotional responses.} \\
\\ 
\texttt{To score Remy's emotional state at the end of the conversation, I will consider the following:} \\
\\ 
\texttt{* Remy's romanticism and willingness to make a wager with Eilis suggests that he is feeling optimistic and enthusiastic.} \\
\texttt{* However, Eilis's rejection and mocking tone may have left Remy feeling disappointed and possibly even hurt.} \\
\texttt{* Remy's final statement, "And you were always afraid to lose, Eilis," suggests that he is trying to assert dominance and challenge Eilis's emotional vulnerability.} \\
\\ 
\texttt{Based on these observations, I will score Remy's emotional state as follows:} \\
\\ 
\texttt{Empathetic: 0} \\
\texttt{Mocking: 0} \\
\texttt{Astonished: 0} \\
\texttt{Disappointed: 8} \\
\\ 
\\ 
\\ 
\texttt{Your task is to predict the likely emotional responses of a character in this dialogue:} \\
\\ 
\texttt{Eilis: Remy, you're late.} \\
\texttt{Remy: And you're as punctual as ever, Eilis. Still chasing shadows, are we?} \\
\texttt{Eilis: Better shadows than illusions of grandeur, Remy.} \\
\texttt{Remy: Ah, the sting of your wit. I've missed that. } \\
\texttt{Eilis: And I've missed your incessant need to prove something.} \\
\texttt{Remy: Is that so? Well, I suppose we're both here for the same reason then.} \\
\texttt{Eilis: Don't flatter yourself, Remy. Our reasons couldn't be more different.} \\
\texttt{Remy: Oh, do enlighten me, Eilis. } \\
\texttt{Eilis: You're here to feed your ego. I'm here to reclaim what I lost.} \\
\texttt{Remy: And what if I told you I'm here for you, Eilis?} \\
\texttt{Eilis: I'd say you're a bigger fool than I thought.} \\
\texttt{Remy: Well, then let's be fools together, shall we? } \\
\texttt{Eilis: I don't need your charity, Remy.} \\
\texttt{Remy: It's not charity, Eilis. It's competition. } \\
\texttt{Eilis: Is that a challenge, Remy?} \\
\texttt{Remy: If you're not afraid, Eilis.} \\
\texttt{Eilis: Afraid? Of you? Don't make me laugh.} \\
\texttt{Remy: Then let's make a wager, Eilis. } \\
\texttt{Eilis: A wager? And what would that be?} \\
\texttt{Remy: The first to create the Philosopher's Stone.} \\
\texttt{Eilis: And what does the loser get, Remy?} \\
\texttt{Remy: The satisfaction of knowing they gave it their all, Eilis.} \\
\texttt{Eilis: You always were a romantic fool, Remy.} \\
\texttt{Remy: And you were always afraid to lose, Eilis.} \\
\texttt{Eilis: I'm not afraid to lose, Remy. I'm afraid to fail.} \\
\texttt{{[}End dialogue{]}} \\
\\ 
\texttt{At the end of this dialogue, Remy would feel...} \\
\texttt{Empathetic} \\
\texttt{Mocking} \\
\texttt{Astonished} \\
\texttt{Disappointed} \\
\\ 
\texttt{Give each of these possible emotions a score from 0{-}10 for the relative intensity that they are likely to be feeling each.} \\
\\ 
\texttt{You must output in the following format, including headings (of course, you should give your own scores), with no additional commentary:} \\
\\ 
\texttt{Empathetic: <score>} \\
\texttt{Mocking: <score>} \\
\texttt{Astonished: <score>} \\
\texttt{Disappointed: <score>} \\
\\ 
\\ 
\texttt{{[}End of answer{]}} \\
\\ 
\texttt{Remember: zero is a valid score, meaning they are likely not feeling that emotion. You must score at least one emotion > 0.} \\
\\ 
\texttt{Your answer:} \\
            }
        }
    }
     \\ \\

    \theutterance \stepcounter{utterance}  
    & \multicolumn{4}{p{0.6\linewidth}}{
        \cellcolor[rgb]{0.561,0.737,0.561}{
            \makecell[{{p{\linewidth}}}]{
                \texttt{\tiny{[P1$\rangle$GM]}}
                \texttt{Empathetic: 2} \\
\texttt{Mocking: 8} \\
\texttt{Astonished: 0} \\
\texttt{Disappointed: 9} \\
            }
        }
    }
    & & \\ \\

    \theutterance \stepcounter{utterance}  
    & & & \multicolumn{2}{p{0.3\linewidth}}{
        \cellcolor[rgb]{0.9,0.9,0.9}{
            \makecell[{{p{\linewidth}}}]{
                \texttt{\tiny{[GM$|$GM]}}
                \texttt{Empathetic: 2} \\
\texttt{Mocking: 8} \\
\texttt{Astonished: 0} \\
\texttt{Disappointed: 9} \\
            }
        }
    }
    & & \\ \\

    \theutterance \stepcounter{utterance}  
    & & & \multicolumn{2}{p{0.3\linewidth}}{
        \cellcolor[rgb]{0.9,0.9,0.9}{
            \makecell[{{p{\linewidth}}}]{
                \texttt{\tiny{[GM$|$GM]}}
                \texttt{Empathetic: 6} \\
\texttt{Mocking: 7} \\
\texttt{Astonished: 0} \\
\texttt{Disappointed: 2} \\
            }
        }
    }
    & & \\ \\

    \theutterance \stepcounter{utterance}  
    & & & \multicolumn{2}{p{0.3\linewidth}}{
        \cellcolor[rgb]{0.9,0.9,0.9}{
            \makecell[{{p{\linewidth}}}]{
                \texttt{\tiny{[GM$|$GM]}}
                \texttt{game\_result = LOSE} \\
            }
        }
    }
    & & \\ \\

    \theutterance \stepcounter{utterance}  
    & & \multicolumn{4}{p{0.6\linewidth}}{
        \cellcolor[rgb]{0.95,0.95,0.95}{
            \makecell[{{p{\linewidth}}}]{
                \texttt{\tiny{[P1$\leftrightarrow$P]}}
                \texttt{Post{-}episode Reflexion (not sent to GM):} \\
\texttt{New plan: } \\
\\ 
\texttt{Upon re{-}reading the dialogue, I realized that I focused too much on the surface{-}level banter between Remy and Eilis, without delving deeper into their underlying emotions and motivations. I should have paid more attention to the subtle cues and power dynamics at play in their conversation.} \\
\\ 
\texttt{To improve my analysis, I will focus on identifying the emotional triggers and vulnerabilities that Remy and Eilis expose throughout the dialogue. I will also consider the context of their relationship and the stakes of their wager, which may have influenced their emotional responses.} \\
\\ 
\texttt{Specifically, I will:} \\
\\ 
\texttt{1. Identify the moments where Remy and Eilis reveal their emotional vulnerabilities, such as Eilis's fear of failure and Remy's romanticism.} \\
\texttt{2. Analyze the power dynamics at play in their conversation, including the ways in which they use language to assert dominance or challenge each other.} \\
\texttt{3. Consider the context of their relationship and the implications of their wager, which may have influenced their emotional responses.} \\
\\ 
\texttt{To score Remy's emotional state at the end of the conversation, I will consider the following:} \\
\\ 
\texttt{* Remy's romanticism and willingness to make a wager with Eilis suggests that he is feeling optimistic and enthusiastic.} \\
\texttt{* However, Eilis's rejection and mocking tone may have left Remy feeling disappointed and possibly even hurt.} \\
\texttt{* Remy's final statement, "And you were always afraid to lose, Eilis," suggests that he is trying to assert dominance and challenge Eilis's emotional vulnerability.} \\
\\ 
\texttt{Based on these observations, I will score Remy's emotional state as follows:} \\
\\ 
\texttt{Empathetic: 0} \\
\texttt{Mocking: 0} \\
\texttt{Astonished: 0} \\
\texttt{Disappointed: 8} \\
            }
        }
    }
    & \\ \\

\end{supertabular}
}

\end{document}
